% ═══════════════════════════════════════════════════════════════════════════════
%   Kuratowski Calculus Engine — Unified Code Manual
%   XeLaTeX document · Styled after Modulo Synthesis Vol II
% ═══════════════════════════════════════════════════════════════════════════════
%
%   Compile with:
%     xelatex causal_set_sim.tex && xelatex causal_set_sim.tex
%
\documentclass[11pt, a4paper, openany]{book}

% ─── Geometry ─────────────────────────────────────────────────────────────────
\usepackage[
  top=2.5cm, bottom=2.5cm,
  left=3cm, right=3cm,
  headheight=14pt
]{geometry}

% ─── Fonts (XeLaTeX) ─────────────────────────────────────────────────────────
\usepackage{fontspec}
\setmainfont{LMRoman10}[
  BoldFont=LMRoman10-Bold,
  ItalicFont=LMRoman10-Italic,
  BoldItalicFont=LMRoman10-BoldItalic
]
\setmonofont{LMMonoLt10}[
  BoldFont=LMMonoLt10-Bold
]

% ─── Colors (Modulo Synthesis palette) ────────────────────────────────────────
\usepackage{xcolor}
\definecolor{structure}{HTML}{1A535C}   % Deep Teal
\definecolor{main}{HTML}{C6892A}        % Golden Ochre
\definecolor{second}{HTML}{2E4057}      % Slate Blue
\definecolor{third}{HTML}{8B2635}       % Burgundy
\definecolor{fourth}{HTML}{3A7D44}      % Forest Green
\definecolor{mainlight}{HTML}{FDF6E3}   % Solarized Base3
\definecolor{secondlight}{HTML}{EDF2F4}
\definecolor{thirdlight}{HTML}{F9EBEC}
\definecolor{fourthlight}{HTML}{EBF5EC}

% ─── Packages ─────────────────────────────────────────────────────────────────
\usepackage{amsmath,amssymb,amsthm,mathtools}
\usepackage{booktabs,enumitem,microtype,float,longtable,array}
\usepackage{fancyhdr,titlesec,titletoc}
\usepackage{tcolorbox}
\tcbuselibrary{skins,breakable,theorems}
\usepackage{etoolbox}
\usepackage{hyperref}
\hypersetup{
  colorlinks=true,
  linkcolor=structure, citecolor=second, urlcolor=second,
  bookmarksnumbered=true, pdfstartview=FitH,
  pdftitle={Kuratowski Calculus Engine — Code Manual},
  pdfauthor={Juan Pablo Silva Alvarado}
}

% ─── Macros ───────────────────────────────────────────────────────────────────
\newcommand{\lP}{\ell_{\mathrm{P}}}
\newcommand{\dS}{d_{\mathrm{S}}}
\newcommand{\MPl}{M_{\mathrm{Pl}}}
\renewcommand{\chaptername}{Chapter}
\newcommand{\code}[1]{\texttt{#1}}
\newcommand{\func}[1]{\texttt{\color{structure}#1}}
\newcommand{\modname}[1]{\texttt{\bfseries\color{main}#1}}

% ─── Chapter/Section Styling ─────────────────────────────────────────────────
\titleformat{\chapter}[display]
  {\normalfont\Large\bfseries\color{structure}}
  {\small\color{structure!60}\MakeUppercase{\chaptertitlename}\ \thechapter}
  {4pt}{\Huge}
\titlespacing*{\chapter}{0pt}{-20pt}{20pt}

\titleformat{\section}
  {\normalfont\Large\bfseries\color{structure}}{\thesection}{0.6em}{}
\titlespacing*{\section}{0pt}{12pt}{6pt}

\titleformat{\subsection}
  {\normalfont\large\bfseries\color{structure!80!black}}{\thesubsection}{0.5em}{}
\titlespacing*{\subsection}{0pt}{8pt}{4pt}

% ─── Headers ──────────────────────────────────────────────────────────────────
\pagestyle{fancy}
\fancyhf{}
\fancyhead[LE,RO]{\small\thepage}
\fancyhead[LO]{\small\nouppercase{\rightmark}}
\fancyhead[RE]{\small\nouppercase{\leftmark}}
\renewcommand{\headrulewidth}{0.3pt}

% ─── Theorem Environments (from Vol II) ──────────────────────────────────────
\newcounter{thm}[chapter]
\renewcommand{\thethm}{\thechapter.\arabic{thm}}

% Definition Box
\newtcolorbox[use counter=thm]{definition}[1][]{%
  enhanced,breakable,
  colback=mainlight,colframe=main,
  colbacktitle=main,coltitle=white,
  fonttitle=\small\bfseries,
  title={Definition~\thethm\ifx&#1&\else: #1\fi},
  arc=1.5pt,boxrule=0.6pt,
  left=8pt,right=8pt,top=6pt,bottom=6pt,
  before skip=10pt,after skip=10pt
}

% Theorem Box
\newtcolorbox[use counter=thm]{theorem}[1][]{%
  enhanced,breakable,
  colback=secondlight,colframe=second,
  colbacktitle=second,coltitle=white,
  fonttitle=\small\bfseries,
  title={Theorem~\thethm\ifx&#1&\else: #1\fi},
  arc=1.5pt,boxrule=0.6pt,
  left=8pt,right=8pt,top=6pt,bottom=6pt,
  before skip=10pt,after skip=10pt
}

% Axiom Box
\newtcolorbox[use counter=thm]{axiom}[1][]{%
  enhanced,breakable,
  colback=thirdlight,colframe=third,
  colbacktitle=third,coltitle=white,
  fonttitle=\small\bfseries,
  title={Axiom~\thethm\ifx&#1&\else: #1\fi},
  arc=1.5pt,boxrule=0.6pt,
  left=8pt,right=8pt,top=6pt,bottom=6pt,
  before skip=10pt,after skip=10pt
}

% Result Box
\newtcolorbox{result}[1][]{%
  enhanced,breakable,
  colback=fourthlight,colframe=fourth,
  fontupper=\itshape,
  arc=1.5pt,boxrule=0.6pt,
  left=8pt,right=8pt,top=6pt,bottom=6pt,
  before skip=10pt,after skip=10pt,
  title={Key Result},
  coltitle=fourth!50!black, fonttitle=\small\bfseries,
  attach boxed title to top left={yshift=-2mm,xshift=2mm},
  boxed title style={colback=fourthlight,frame hidden},
  #1
}

% Introduction/Overview Box
\newenvironment{introduction}[1][Chapter Overview]{%
  \begin{tcolorbox}[
    enhanced,colback=structure!5,colframe=structure,
    colbacktitle=structure,coltitle=white,
    fonttitle=\small\bfseries,title={#1},
    arc=2pt,boxrule=0.7pt,
    left=8pt,right=8pt,top=6pt,bottom=6pt,
    before skip=12pt,after skip=12pt
  ]\itshape\begin{itemize}[leftmargin=*,itemsep=2pt,parsep=0pt]
}{%
  \end{itemize}\end{tcolorbox}
}

\setlist{nosep,leftmargin=*}
\setlist[itemize,1]{label={\color{structure}\textbullet}}
\setlist[enumerate,1]{label={\color{structure}\arabic*.}}

% ─── Title Page Macro (from Vol II) ──────────────────────────────────────────
\makeatletter
\newcommand{\subtitle}[1]{\gdef\@subtitle{#1}}
\newcommand{\extrainfo}[1]{\gdef\@extrainfo{#1}}
\gdef\@subtitle{}\gdef\@extrainfo{}
\renewcommand{\maketitle}{%
  \begin{titlepage}\centering
    \vspace*{4cm}
    {\Huge\bfseries\color{structure}\@title\par}\vspace{0.6cm}
    \ifx\@subtitle\empty\else{\Large\color{gray!80}\@subtitle\par}\fi
    \vspace{2cm}
    {\color{main}\rule{0.6\textwidth}{2pt}\par}\vspace{2cm}
    {\LARGE\@author\par}\vspace{1cm}
    {\large\@date\par}\vfill
    \ifx\@extrainfo\empty\else
      \begin{minipage}{0.8\textwidth}\centering\itshape\color{gray!60}\@extrainfo\end{minipage}\vspace{3cm}
    \fi
  \end{titlepage}
}
\makeatother

% ═══════════════════════════════════════════════════════════════════════════════
\title{Kuratowski Calculus Engine}
\subtitle{Unified Code Manual}
\author{Juan Pablo Silva Alvarado}
\date{\today}
\extrainfo{%
  Part of Modulo Synthesis / Fractal Entropic Geometrodynamics\\[4pt]
  ``The universe is strictly computable, finite, and structurally homeostatic.
  The rest is just counting.''}

\begin{document}
\maketitle
\frontmatter
\tableofcontents
\mainmatter

% ███████████████████████████████████████████████████████████████████████████████
\chapter{Theory}
% ███████████████████████████████████████████████████████████████████████████████

\begin{introduction}
  \item The two axioms of Fractal Entropic Geometrodynamics.
  \item The Sole Principle of Interaction (microscopic law).
  \item Kuratowski's Theorem and the origin of matter, forces, and generations.
  \item The Three Geometries: Order, Matter, and Logic.
  \item How each physics concept maps to a code module and data structure.
\end{introduction}

% ─────────────────────────────────────────────────────────────────────────────
\section{Axiom 1 --- Causal Sets}

Spacetime is not a continuous manifold but a \textbf{discrete partial order}
--- a directed acyclic graph (DAG) where each node is a spacetime event and
each edge encodes causal influence: $A \to B$ means ``$A$ can causally
influence~$B$.''  The DAG is generated by Poisson sprinkling $N$ points into a
4D causal diamond ($|t| + r \le T/2$) with fundamental density $\rho \approx 1$.

\begin{axiom}[Causal Sets (from Vol~I)]
Spacetime is a locally finite partially ordered set.  The order relation
encodes causal structure; the counting measure (number of elements in an
Alexandrov set) encodes spacetime volume.  All geometry is emergent from
order + number.
\end{axiom}

% ─────────────────────────────────────────────────────────────────────────────
\section{Axiom 2 --- Quantum Dynamics}

\begin{axiom}[Quantum Dynamics: The Non-Commutative Postulate (from Vol~II)]
The geometry and quantum phases of the causal set are evaluated over a
\textbf{non-commutative von Neumann algebra}.  Macroscopic geometry is not
fundamental; it is an emergent phenomenon dictated by the quantum Fisher
Information Metric (Petz metric) defined over the density states of causal
histories.  Causal histories do not commute: if events $A$ and $B$ are
spacelike separated, the order in which they are evaluated affects the quantum
phase.
\end{axiom}

Together, Chentsov's Theorem (classical, macroscopic) and Petz's Theorem
(quantum, microscopic) establish a single information-geometric foundation for
all of physics.  Gravity is the coarse-grained Fisher geometry at large~$N$;
quantum mechanics is the fine-grained Petz geometry at small~$N$.

% ─────────────────────────────────────────────────────────────────────────────
\section{The Sole Principle of Interaction}

\begin{axiom}[The Sole Principle of Interaction (Microscopic Law)]
The entirety of microscopic network dynamics is governed by a single mandate:
\textbf{the preservation of unitarity} (probability conservation) across
discrete, non-commutative causal links:
\[
  \sum_{\text{paths}} P(i \to f) = 1
\]
All emergent forces (gauge fields) and geometric structures (gravity) are
\textbf{homeostatic mechanisms} enacted by the causal network to satisfy this
unitary principle under the topological constraints of Kuratowski's Theorem.
\end{axiom}

% ─────────────────────────────────────────────────────────────────────────────
\section{Hasse Diagram}

The \textbf{transitive reduction} of the causal DAG.  If $A \to B \to C$, the
direct edge $A \to C$ is redundant and is removed.  The Hasse diagram retains
only nearest-neighbour causal links --- the skeleton of spacetime.

% ─────────────────────────────────────────────────────────────────────────────
\section{Causal Prism --- The Fundamental Particle}

\textbf{CRITICAL:} The transitive reduction in Phase 1 produces a
\emph{triangle-free} DAG.  $K_4$ complete subgraphs (4 mutually connected
events) \textbf{cannot exist} in such a graph --- transitive reduction
eliminates all triangles by definition.  The correct fundamental
particle is the \textbf{Causal Prism} ($K_{2,N}$ bipartite, $N \ge 3$).

A Causal Prism consists of:
\begin{itemize}
  \item Two \textbf{poles}: past pole $u$ (origin) and future pole $v$ (destination)
  \item $N \ge 3$ \textbf{belly nodes} (intermediates $w_i$)
  \item Directed edges: $u \to w_i$ and $w_i \to v$ for all $i \in \{1 \dots N\}$
  \item The belly = $\text{children}(u) \cap \text{parents}(v)$
\end{itemize}

These topological defects:
\begin{itemize}
  \item Anchor probability flow and define particle-like spectral signatures.
  \item Cannot be ``untied'' without breaking the graph's connectivity ---
        they are topologically stable.
  \item Are the \emph{maximal} bipartite structures permitted by the
        triangle-free constraint.
\end{itemize}

% ─────────────────────────────────────────────────────────────────────────────
\section{Threat Contraction --- Colour Confinement}

\begin{theorem}[Kuratowski's Theorem (The Topological UV Limit)]
A finite graph $G$ is planar if and only if it does not contain a topological
minor of $K_5$ (the complete graph on 5 vertices) or $K_{3,3}$ (the complete
bipartite graph on $3+3$ vertices):
\[
  K_5 \not\preceq_T G_{\text{UV}} \quad \text{and} \quad
  K_{3,3} \not\preceq_T G_{\text{UV}}
\]
This is the absolute UV boundary of the theory.
\end{theorem}

When an external node is connected to \textbf{both poles} of a Causal Prism
AND $\ge 2$ intermediates, it threatens the topological stability and is
\textbf{absorbed} (vertex contraction) into the higher-degree pole --- a
discrete analogue of colour confinement.  The topology enforces that quarks
cannot exist in isolation.

\begin{result}
The Strong Force is not a fundamental gauge field.  It is a strict
topological obstruction --- the causal network's response to the Kuratowski UV
limit --- preventing forbidden topological minors in the ultraviolet.
\end{result}

% ─────────────────────────────────────────────────────────────────────────────
\section{Spectral Dimension --- Gravity}

The spectral dimension measures how quickly a random walker returns to its
starting point after $t$ steps:
\[
  \dS(t) = -2\,\frac{d(\ln P)}{d(\ln t)}
\]
On a smooth 4D manifold, $\dS \to 4$.  On the Hasse graph:
\begin{itemize}
  \item \textbf{Vacuum:} $\dS$ diverges (boundary saturation of the
        unstructured bulk).
  \item \textbf{With Causal Prism matter:} $\dS$ stabilises at a finite value
        --- the topological knot traps probability flow.
\end{itemize}

% ─────────────────────────────────────────────────────────────────────────────
\section{Causal Flux --- Electromagnetism}

\textbf{Directed walkers} follow the arrow of time (cause $\to$ effect only).
The transmission probability between particle generations measures
electromagnetic-like interactions:
\begin{itemize}
  \item \textbf{Attraction:} Gen1 $\to$ AntiGen1 (opposite charges).
  \item \textbf{Repulsion:} Gen1 $\to$ Gen1 (same charges).
\end{itemize}

% ─────────────────────────────────────────────────────────────────────────────
\section{Three Generations --- The Kuratowski Limit}

The Standard Model has three copies of matter (electron/muon/tau,
up/charm/top).  In FEG the answer is a strict topological obstruction:
\begin{enumerate}
  \item UV gravity compresses $\dS \to 2$ (effectively 2D topology).
  \item Kuratowski's Theorem forbids $K_5$ in planar graphs.
  \item The transitive reduction produces a triangle-free DAG.
  \item Causal Prisms ($K_{2,N}$ bipartite) are the fundamental particles.
\end{enumerate}

Causal Prisms are classified by their bit-packed topological signature
$[q_0, q_1, q_2, q_3]$ (quartile bulk momentum across all components: origin,
destination, and intermediates).  Most abundant signature $=$ Gen1
(electron-like), second $=$ Gen2 (muon-like), third $=$ Gen3 (tau-like).  The
CPT conjugate (reversed and negated components) $=$ AntiGen1 (positron-like).
The emergent discreteness of signatures enforces exactly 3 stable generations.

% ─────────────────────────────────────────────────────────────────────────────
\section{The Three Geometries}

The theory rests on three geometric pillars, each discovered through the
simulation engine:

\begin{result}
\textbf{The Three Geometries of Fractal Entropic Geometrodynamics:}
\begin{enumerate}
  \item \textbf{Geometry of Order} (Vol~I): Transitive Reduction $\to$ Lorentz
        invariance.  The Hasse diagram encodes the causal structure of spacetime;
        counting elements in an Alexandrov set reproduces the volume element.
        \emph{Code:} \code{diamond.rs} builds the Hasse DAG.
  \item \textbf{Geometry of Matter} (Vol~II): Causal Prism $K_{2,N}$ $\to$ mass
        hierarchy.  Topological defects in the Hasse graph carry quantized
        topological mass $N$ and are classified into exactly three generations by
        the Kuratowski limit.
        \emph{Code:} \code{skyrmion.rs} detects prisms and classifies generations.
  \item \textbf{Geometry of Logic}: $O(N)$ scaling $\Leftrightarrow$ physical
        locality.  Bounded Hasse degree ($D \le 15$) implies that every
        topological operation---prism detection, threat contraction, spectral
        walk---is local, yielding \textbf{provably linear} $O(N)$ computational
        complexity.  This is not an engineering optimisation but a
        \emph{computational proof} that the theory encodes physical locality as an
        information principle.
        \emph{Code:} 2-hop candidate search in \code{skyrmion.rs};
        single-pass streaming in \code{diamond.rs}.
\end{enumerate}
\end{result}

\begin{theorem}[$O(N)$ Locality Theorem]
Let $\mathcal{H}$ be the Hasse diagram of $N$ events sprinkled into a 4D causal
diamond at density $\rho \approx 1$.  Then:
\begin{enumerate}
  \item Every Hasse link has bounded proper time: $\tau_{\text{link}} \le 8$.
  \item Hence every node has bounded Hasse degree: $D \le D_{\max} \approx 15$.
  \item Causal Prism detection via 2-hop traversal is $O(N \cdot D_{\max}^3) = O(N)$.
  \item All other phases (edge generation, spectral walkers, output) are already
        $O(N)$ or $O(W \cdot t)$.
\end{enumerate}
The entire simulation scales linearly with the number of spacetime events.
\end{theorem}

% ─────────────────────────────────────────────────────────────────────────────
\section{Physics $\to$ Code Dictionary}

\begin{table}[H]
\centering
\begin{tabular}{@{} l l l @{}}
\toprule
\textbf{Physics Concept} & \textbf{Code Module} & \textbf{Data Structure} \\
\midrule
Spacetime              & \code{diamond.rs}  & DAG of \code{[f64; 4]} points \\
Causal link            & \code{diamond.rs}  & CSR edge \code{(adj\_head, adj\_data)} \\
Particle (matter)      & \code{skyrmion.rs} & Causal Prism ($K_{2,N}$ bipartite) \\
Strong force           & \code{skyrmion.rs} & Threat contraction \code{merge\_into[i]} \\
Charge (bulk momentum) & \code{diamond.rs}  & \code{bulk\_momentum: Vec<i32>} \\
Particle generation    & \code{skyrmion.rs} & Bit-packed signature \code{u32} \\
Antimatter             & \code{skyrmion.rs} & CPT-conjugate signature \\
Gravity (geometry)     & \code{spectral.rs} & Undirected walkers $\to P(t) \to \dS(t)$ \\
Electromagnetism       & \code{spectral.rs} & Directed walkers $\to$ transmission prob. \\
Locality ($O(N)$)      & \code{skyrmion.rs} & 2-hop traversal, bounded degree $D \le 15$ \\
\bottomrule
\end{tabular}
\caption{Mapping between physics concepts and code implementations.}
\label{tab:dictionary}
\end{table}


% ███████████████████████████████████████████████████████████████████████████████
\chapter{Architecture}
% ███████████████████████████████████████████████████████████████████████████████

\begin{introduction}
  \item The four-phase pipeline: sprinkle $\to$ hasse $\to$ defect $\to$ walkers $\to$ csv.
  \item Two execution paths: in-memory vs streaming (single-pass $O(N)$).
  \item Tier system: eigendecomposition vs Monte Carlo.
\end{introduction}

% ─────────────────────────────────────────────────────────────────────────────
\section{Data Flow}

\begin{verbatim}
  sprinkle(N) -> [f64;4] points
      |
  build_hasse() -> CSR adjacency (adj_head, adj_data) + bulk_momentum
      |
  apply_defect() -> Causal Prisms, K5 threat absorption, generation classification
      |
  run_walkers() -> P(t), dS(t) for vacuum/defect/generations/flux
      |
  write_csv() -> results.csv
\end{verbatim}

% ─────────────────────────────────────────────────────────────────────────────
\section{Phase Diagram}

\begin{table}[H]
\centering
\begin{tabular}{@{} c l l @{}}
\toprule
\textbf{Phase} & \textbf{Module} & \textbf{Physics} \\
\midrule
1 & \code{diamond.rs}  & Vacuum generation: Poisson sprinkling + Hasse diagram \\
2 & \code{skyrmion.rs} & Kuratowski contraction: Causal Prism detection + $K_5$ absorption \\
3 & \code{spectral.rs} & Spectral dimension: random walk return probability \\
4 & \code{output.rs}   & CSV serialisation of ensemble-averaged observables \\
\bottomrule
\end{tabular}
\caption{The four simulation phases.}
\end{table}

% ─────────────────────────────────────────────────────────────────────────────
\section{Execution Paths}

\subsection{In-Memory}
Selected by \code{-{}-inmemory} flag or auto-selected when RAM is sufficient.
All data structures live in RAM.  Realisations run in parallel via Rayon with an
auto-selected concurrency limit (4 threads, or 2 for $N > 10$M; overridable via
\code{-{}-threads}).  Causal Prism detection uses 2-hop local candidate search:
$O(N \cdot D^3)$ with $D \le 15$.

\subsection{Streaming}
Selected by \code{-{}-stream} flag or auto-selected when RAM is insufficient.
Phase~1 uses \func{scan\_edges\_with\_analysis()} to perform a \textbf{single-pass}
scan: degree counting, core classification, and core-edge collection are merged
into one streaming pass over the edge generator.  Causal Prism detection then
operates on the core-induced CSR subgraph using the same 2-hop local search as
the in-memory path.  RAM usage stays under ${\sim}1.3$\,GB per realisation even
at $N = 100$M.  Concurrency is auto-selected by CPU count (capped at 8;
overridable via \code{-{}-threads}).

Mode selection is automatic based on available system RAM (via \code{sysinfo}),
configurable via \code{causal\_set.toml}, or overridable with CLI flags.

% ─────────────────────────────────────────────────────────────────────────────
\section{Tier System}

\begin{table}[H]
\centering
\begin{tabular}{@{} l l l l @{}}
\toprule
\textbf{Tier} & \textbf{Condition} & \textbf{Algorithm} & \textbf{Complexity} \\
\midrule
Eigendecomp & $N \le 3{,}000$ (default \code{-{}-eigen-cutoff}) & Dense $S = D^{-1/2}AD^{-1/2}$ eigendecomp & $O(N^3)$ \\
Monte Carlo & $N > 3{,}000$ (default \code{-{}-eigen-cutoff}) & Lazy random walkers on CSR & $O(W \cdot t)$ \\
\bottomrule
\end{tabular}
\caption{Spectral dimension computation tiers.}
\end{table}


% ███████████████████████████████████████████████████████████████████████████████
\chapter{Module Reference: \modname{diamond} --- Phase 1}
% ███████████████████████████████████████████████████████████████████████████████

\begin{introduction}[Vacuum Generation]
  \item Poisson sprinkling in a 4D causal diamond.
  \item Hasse diagram construction (transitive reduction).
  \item Two tiers: sparse $A^2$ reduction ($N \le 15$k) and OCI geometric construction ($N > 15$k).
\end{introduction}

% ─────────────────────────────────────────────────────────────────────────────
\section{Functions}

\subsection{\func{sprinkle(n, rng)} $\to$ \code{(Vec<[f64;4]>, f64)}}

Rejection-samples $n$ points inside $|t| + r \le T/2$ with $T = (24N)^{1/4}$
giving density $\rho \approx 1$.  The causal diamond is the intersection of
the future light cone of $(t = -T/2)$ and the past light cone of $(t = +T/2)$.

\subsection{\func{build\_hasse\_sparse(pts)} $\to$ \code{(coords, adj\_head, adj\_data, momentum)}}

Full causal closure followed by sparse $A^2$ reduction.  Edge $(i,j)$ is a
Hasse link iff $A^2[i,j] = 0$ (no 2-hop path exists).  Used for $N \le
15{,}000$.  Parallel over source nodes via Rayon.

\subsection{\func{build\_hasse\_direct(pts)} $\to$ \code{(coords, adj\_head, adj\_data, momentum)}}

Direct geometric construction using the \textbf{Occupied Cell Index (OCI)} --- a
spatial indexing paradigm that replaces naive geometric shell enumeration (96\%
empty cells) with per-time-layer lists of occupied cells, queried via binary
search.  Features:
\begin{itemize}
  \item Points are quantised onto an integer lattice and sorted by cell index
        for cache-coherent traversal (\code{GridPoint} struct).
  \item Per-time-layer lists of occupied cells, sorted by X coordinate.
  \item Binary search for the active X-band within the light cone.
  \item \textbf{Topological cutoffs} (Phase~1 optimisations):
        \begin{itemize}
          \item \code{MAX\_CAUSAL\_DEPTH = 4}: proper-time cutoff.  Hasse links
                beyond $\Delta t > 4$ are exponentially suppressed at $\rho
                \approx 1$ (was 10).
          \item \code{MAX\_HASSE\_DEGREE = 15}: degree cap (Valve~1) --- a node
                with 15+ descendants has fully populated its causal shadow
                (was unbounded).
        \end{itemize}
  \item \textbf{Alexandrov volume short-circuit (Valve~2):} If every causal
        candidate in a $\Delta t$-shell is transitively blocked (non-empty
        Alexandrov set between $u$ and $v$), further shells have strictly
        larger Alexandrov volume and are guaranteed blocked.  Breaks
        immediately when degree $\ge 1$.
  \item \textbf{Adaptive dry-shell termination (Valve~3):} 1 dry
        $\Delta t$-shell (no causal candidates) $+$ $\ge 2$ descendants
        triggers cutoff.
  \item Hollow light-cone shell optimisation: skip deep interior points where
        proper time $> \code{MAX\_CAUSAL\_DEPTH}$ units.
\end{itemize}

\subsection{\func{stream\_edges\_to\_file(n\_total, chunk\_size, path, seed)}}

Streaming variant that writes Hasse edges to disk using a sliding window, for
$N$ too large to fit in RAM.  Output format: continuous stream of
\code{(u32, u32)} in little-endian.  Both physics invariants are preserved: full
light cone search and transitive reduction via locally cached children
coordinates.

\subsection{\func{sprinkle\_chunk(n\_expected, t\_min, t\_max, half\_t, rng, start\_idx)}}

Poisson-sprinkle points into a time-slab $[t_{\min}, t_{\max}]$.  Used
internally by the streaming path to fill the sliding window buffer.

\subsection{\func{scan\_edges\_sparse\_ext(n\_total, chunk\_size, seed, on\_batch)}}

Extended streaming edge scanner.  Identical to \func{scan\_edges\_sparse} except
the callback receives the current \code{safe\_threshold} as a second argument:
\code{FnMut(\&[(u32, u32)], f64)}.  This enables callers to track node
finalization (a node's total degree is known once
$\text{safe\_threshold} > t_{\text{node}} + \text{safety\_margin}$).

\subsection{\func{scan\_edges\_with\_analysis(n\_total, chunk\_size, seed, core\_num, core\_den)}}

\textbf{Single-pass streaming analysis} that merges degree counting, core
classification, and core-edge collection into one scan over the edge generator.
Returns \code{(in\_deg, out\_deg, is\_core, directed\_core\_edges)}.

Pipeline inside the callback:
\begin{enumerate}
  \item Update \code{in\_deg[v]} and \code{out\_deg[u]} for each edge.
  \item Track finalized nodes via safe-threshold delay
        ($\text{safe\_threshold} - 2 \cdot \text{max\_dt} - 4.0$).
  \item Running degree histogram ($10{,}000$ bins) estimates the core threshold
        after 5\% warmup.
  \item Classify finalized nodes as core or non-core.
  \item Collect edges where both endpoints are classified core; buffer edges
        where one endpoint is pending.
  \item Final exact $O(N)$ histogram pass corrects any drift ($\pm 0.5\%$) and
        re-checks all collected core edges.
\end{enumerate}

\begin{result}
The single-pass architecture eliminates the two-pass bottleneck (degree counting
+ core-edge collection), halving wall-clock time for $N > 1$M streaming runs.
At $N = 100$M each realisation uses ${\sim}1.3$\,GB RAM ($13 \cdot N$ bytes).
\end{result}

% ─────────────────────────────────────────────────────────────────────────────
\section{Key Concepts}

\begin{definition}[GridPoint]
A spacetime event quantised onto an integer lattice for cache-coherent
traversal.  Fields: \code{p} (coordinates), \code{orig\_idx}, \code{cell}
(linearised 4D grid coordinate), \code{qt/qx/qy/qz} (quantised coordinates).
Points are sorted by cell index so that spatially nearby events are contiguous
in memory.
\end{definition}

\begin{definition}[Occupied Cell Index (OCI)]
Paradigm shift from geometric to spatial indexing.  The naive approach iterates
all cells in the light-cone shell for each source node.  At large $N$ the grid
is sparse ($>96\%$ cells empty), making geometric enumeration wasteful.  OCI
inverts the lookup:
\begin{enumerate}
  \item Build per-time-layer lists of \emph{occupied} cells only.
  \item Sort each layer by X coordinate.
  \item For each source, binary-search the X-band of the light cone within the
        occupied list, skipping all empty space.
\end{enumerate}
This reduces the inner-loop cost from $O(\text{shell\_volume})$ to
$O(\text{occupied\_cells})$, a ${\sim}20\times$ speedup at $N > 100$k.
\end{definition}


% ███████████████████████████████████████████████████████████████████████████████
\chapter{Module Reference: \modname{skyrmion} --- Phase 2}
% ███████████████████████████████████████████████████████████████████████████████

\begin{introduction}[Kuratowski Contraction]
  \item Pure integer topology --- zero floating-point.
  \item Core identification $\to$ Causal Prism ($K_{2,N}$ bipartite) detection $\to$ threat contraction $\to$ generation classification.
  \item Bit-packed topological signature encoding (quartile bulk momentum).
\end{introduction}

% ─────────────────────────────────────────────────────────────────────────────
\section{Algorithm}

\textbf{CRITICAL:} The transitive reduction produces a triangle-free DAG.
$K_4$ complete subgraphs \textbf{cannot exist}.  The fundamental particle is
the \textbf{Causal Prism} ($K_{2,N}$ bipartite, $N \ge 3$): two poles with
$N \ge 3$ belly nodes (intermediates satisfying $u \to w_i \to v$).

\begin{enumerate}
  \item \textbf{Core selection:} Top 10\% of nodes by undirected Hasse degree
        ($\approx$ the diamond's combinatorial centre).
  \item \textbf{Causal Prism formation ($O(N)$ forward-forward 2-hop search):}
        \textbf{Minimum valence pruning:} nodes with out-degree
        $< \code{MIN\_PRISM\_SHARED}$ (3) are skipped before the 2-hop
        traversal, and candidate poles~$v$ with in-degree $< 3$ are skipped
        before the belly intersection.  For each remaining core node~$u$,
        collect 2-hop candidates by traversing $u \to w \to v$ through forward
        successors~$w$ (children of~$u$) and their forward successors~$v$
        (grandchildren of~$u$, future pole candidates).  Then count belly
        nodes $|\text{children}(u) \cap \text{parents}(v)|$ using the reverse
        CSR for $\text{parents}(v)$.  If $\ge 3$, the pair forms a Prism
        candidate with the belly as intermediates.  Greedy best-partner
        selection maximizes belly size per pole.

        \textbf{Complexity:} $O(N_{\text{core}} \cdot D^3)$.  With bounded Hasse
        degree $D \le 15$: $O(10\text{M} \times 3{,}375) \approx 3.4 \times 10^{10}$
        operations at $N = 100$M --- minutes instead of days.

        \textbf{Completeness proof:} If $u$ and $v$ are poles of a timelike
        $K_{2,N}$, then every belly node~$w_i$ satisfies $u \to w_i \to v$.
        The forward-forward path $u \to w_i \to v$ guarantees $v$ appears as a
        candidate of~$u$.  Every valid prism partner \emph{must} have at least
        3 belly nodes, each providing such a forward-forward path.  The 2-hop
        search is therefore \textbf{complete}: no valid partner is missed.
        The valence guard is safe because any valid pole must have $\ge 3$
        children (belly nodes).
  \item \textbf{Threat contraction:} Any external node connected to BOTH poles
        AND $\ge 2$ intermediates is absorbed into the higher-degree pole (vertex
        contraction).
  \item \textbf{Generation classification:} Prisms classified by bit-packed
        topological signature (quartile bulk momentum across all components).
  \item \textbf{Transitive merge chain resolution:} Pointer chase to resolve
        chains ($O(N)$).
  \item \textbf{Defect CSR construction:} Zero-copy CSR from merge map +
        pole-completion edges (one edge per Prism: origin $\leftrightarrow$ destination).
\end{enumerate}

% ─────────────────────────────────────────────────────────────────────────────
\section{Functions}

\subsection{\func{apply\_defect(n, adj\_head, adj\_data, bulk\_momentum)} $\to$ \code{DefectResult}}

Main entry point.  Finds Causal Prisms ($K_{2,N}$ bipartite structures with
belly nodes), classifies them by topological signature into generations,
performs threat contraction, and builds the defect CSR by adding
pole-completion edges.

\subsection{\func{process\_streaming(n\_total, chunk\_size, path, seed)} $\to$ \code{SpectralResult}}

Out-of-core variant for streaming mode.  Single-pass analysis via
\func{scan\_edges\_with\_analysis()}:
\begin{enumerate}
  \item \textbf{Single pass:} Stream edges through
        \func{scan\_edges\_with\_analysis()} to simultaneously count degrees,
        classify core nodes, and collect core-induced edges.  Returns
        \code{(in\_deg, out\_deg, is\_core, directed\_core\_edges)}.
  \item \textbf{Prism detection:} Build CSR for the core-induced subgraph, run
        $O(N)$ 2-hop Causal Prism detection and generation classification.
  \item \textbf{Spectral measurement:} Run random walkers on the core CSR for
        spectral dimension; build directed CSR for causal flux.
\end{enumerate}

% ─────────────────────────────────────────────────────────────────────────────
\section{\code{DefectResult} Struct}

\begin{table}[H]
\centering
\begin{tabular}{@{} l p{9cm} @{}}
\toprule
\textbf{Field} & \textbf{Description} \\
\midrule
\code{vac\_head}, \code{vac\_data}   & Vacuum Hasse CSR (row pointers + column indices). \\
\code{def\_head}, \code{def\_data}   & Defect CSR after threat contraction + Prism pole-completion edges. \\
\code{vacuum\_core}                  & Core node indices (top 10\% by degree) in the vacuum graph. \\
\code{defect\_core}                  & Core node indices surviving threat contraction. \\
\code{gen1\_nodes}                   & Generation 1 Prism node indices --- most abundant signature (electron-like). \\
\code{gen2\_nodes}                   & Generation 2 Prism node indices --- second most abundant (muon-like). \\
\code{gen3\_nodes}                   & Generation 3 Prism node indices --- third most abundant (tau-like). \\
\code{anti1\_nodes}                  & Anti-Generation 1 Prism node indices --- CPT conjugate (positron-like). \\
\code{sterile\_nodes}                & Sterile prism node indices --- prisms with $N > 5$ intermediates (dark matter candidates, Conjecture~C6). \\
\code{merge\_into}                   & Merge map: \code{merge\_into[i]} = canonical node for node~$i$. For threat-absorbed nodes, points to the absorbing Prism pole. \\
\bottomrule
\end{tabular}
\caption{\code{DefectResult} fields.}
\end{table}

% ─────────────────────────────────────────────────────────────────────────────
\section{Signature Encoding}

Each Causal Prism signature is computed from the bulk momentum of ALL
components (origin, destination, and intermediates):
\begin{enumerate}
  \item Collect bulk momentum $[M_{\text{origin}}, M_{\text{destination}},
        M_{w_1}, M_{w_2}, \dots]$ for all nodes in the Prism.
  \item Sort the list and extract quartile values: $[q_0, q_1, q_2, q_3]$
        (0th, 25th, 50th, 75th percentiles).
  \item Normalise each quartile to $[-10, +10]$ via $q_i \times 10 / \max(|q|)$.
  \item Offset by $+16$ to make non-negative (range $[6, 26]$, fits in 5 bits).
  \item Pack four 5-bit components into a single \code{u32}:
\end{enumerate}
\[
  \text{sig} = c_0 \mathbin{|} (c_1 \ll 5) \mathbin{|} (c_2 \ll 10) \mathbin{|} (c_3 \ll 15)
\]
The anti-signature (CPT conjugate) is computed by reversing and negating:
$\bar{c}_i = -c_{3-i} + 16$.


% ███████████████████████████████████████████████████████████████████████████████
\chapter{Module Reference: \modname{spectral} --- Phase 3}
% ███████████████████████████████████████████████████████████████████████████████

\begin{introduction}[Spectral Dimension Measurement]
  \item Random walk return probability via undirected and directed walkers.
  \item Two tiers: eigendecomposition (exact, $N \le 3$k by default; adjustable via \code{-{}-eigen-cutoff}) and Monte Carlo (above the cutoff).
  \item Per-generation and causal flux measurements.
\end{introduction}

% ─────────────────────────────────────────────────────────────────────────────
\section{Functions}

\subsection{\func{spectral\_dimension(steps, p\_vals)} $\to$ \code{Vec<f64>}}

$\dS(t) = -2\,d(\ln P)/d(\ln t)$ via centred finite differences.  Forward
difference at $t=1$, backward at $t=t_{\max}$, centred elsewhere.

\subsection{\func{run\_walkers(adj\_head, adj\_data, starts, steps, seed)} $\to$ \code{Vec<f64>}}

Run $W$ independent lazy-walk walkers in parallel.  At each step: 50\% stay,
50\% move to a uniformly random neighbour.  Returns $P(t)$ = fraction of
walkers that returned to their starting node at each step.  Parallel via Rayon.

\textbf{Strict Finitism:} all internal accumulation uses \code{u64} integer
arithmetic.  Each walker contributes exactly 0 or 1 (Kronecker delta) per
measurement step.  The reduction sums \code{u64} counts across all walkers.
The single \code{f64} division \code{count / W} occurs only at the final
return, preserving exact integer semantics throughout the Monte Carlo core.

\subsection{\func{run\_transmission\_walkers(\ldots)} $\to$ \code{(Vec<f64>, Vec<f64>)}}

Directed walkers (mandatory move along causal arrow).  Returns two vectors:
attraction flux (Gen1 $\to$ AntiGen1) and repulsion flux (Gen1 $\to$ Gen1).
Walkers that reach a dead end (edge of causal set) escape and stop.  Same
\code{u64} integer accumulation as \func{run\_walkers}.

\subsection{\func{compute\_eigen(n, rows, cols, steps, core)} $\to$ \code{SpectralResult}}

Exact $P(t) = N^{-1} \sum_k \lambda_k^t$ from eigendecomposition of the
symmetric normalised adjacency $S = D^{-1/2}AD^{-1/2}$.  Complexity: $O(N^3)$.
Used only for small graphs ($N \le 3$k).  Local $P(t)$ uses
eigenvector-weighted sums over core indices.

\subsection{\func{compute\_monte\_carlo(n, rows, cols, steps, core, n\_w, rng)} $\to$ \code{SpectralResult}}

Monte Carlo $P(t)$ from random walkers.  Builds CSR adjacency internally, then
launches lazy random walkers.  Complexity: $O(W \cdot t)$, independent of~$N$.

\subsection{\func{compute\_monte\_carlo\_csr(n, head, data, steps, core, n\_w, rng)} $\to$ \code{SpectralResult}}

Zero-copy variant --- avoids rebuilding adjacency when CSR is already available
from Phase~2.  This is the hot path for large~$N$ in-memory runs.

\subsection{\func{build\_adj\_list(n, rows, cols)} $\to$ \code{(Vec<u32>, Vec<u32>)}}

Build sorted undirected adjacency list in CSR format from edge lists.

% ─────────────────────────────────────────────────────────────────────────────
\section{\code{SpectralResult} Struct}

\begin{table}[H]
\centering
\begin{tabular}{@{} l p{9cm} @{}}
\toprule
\textbf{Field} & \textbf{Description} \\
\midrule
\code{p\_global}         & $P(t)$ from uniformly random starts. Probes global geometry. \\
\code{ds\_global}        & $\dS(t)$ of the bulk (derived from \code{p\_global}). \\
\code{p\_local}          & $P(t)$ from walkers starting on core nodes. \\
\code{ds\_local}         & $\dS(t)$ from core walkers. Sensitive to topological defects. \\
\code{p\_gen1}, \code{ds\_gen1}   & $P(t)$ and $\dS$ for Generation~1 Causal Prism nodes (electron-like). \\
\code{p\_gen2}, \code{ds\_gen2}   & $P(t)$ and $\dS$ for Generation~2 Prism nodes (muon-like). \\
\code{p\_gen3}, \code{ds\_gen3}   & $P(t)$ and $\dS$ for Generation~3 Prism nodes (tau-like). \\
\code{p\_anti1}, \code{ds\_anti1} & $P(t)$ and $\dS$ for Anti-Generation~1 (positron-like, CPT conjugate). \\
\code{flux\_attraction}  & Directed transmission: Gen1 $\to$ AntiGen1 (opposite charge). \\
\code{flux\_repulsion}   & Directed transmission: Gen1 $\to$ Gen1 (same charge). \\
\code{flux\_attr\_norm}  & Normalized attraction flux: \code{flux\_attraction} / $|\text{targets}|$. Isolates intrinsic per-node coupling strength (Conjecture~C3). \\
\code{flux\_repu\_norm}  & Normalized repulsion flux: \code{flux\_repulsion} / $|\text{targets}|$. \\
\code{p\_sterile}, \code{ds\_sterile} & $P(t)$ and $\dS$ for walkers on sterile prism nodes ($N > 5$). Conjecture~C6: gravitationally active, electromagnetically silent (dark matter candidates). \\
\code{ds\_*\_std}        & Population standard deviation of each $\dS(t)$ field across $M$ realisations. Empty when $M = 1$. Provides error bars for publication figures. \\
\code{flux\_*\_std}      & Population standard deviation of each flux field across $M$ realisations. \\
\bottomrule
\end{tabular}
\caption{\code{SpectralResult} fields (28 total).}
\end{table}


% ███████████████████████████████████████████████████████████████████████████████
\chapter{Module Reference: \modname{output} --- Phase 4}
% ███████████████████████████████████████████████████████████████████████████████

\begin{introduction}[CSV Serialisation]
  \item Ensemble-averaged observables written to CSV.
  \item 32-column format: return probabilities, spectral dimensions, raw and normalized
        causal flux, sterile prism observables, mass spectrum, and ensemble standard
        deviations (error bars for $M > 1$ runs).
\end{introduction}

\section{Function}

\subsection{\func{write\_csv(path, steps, vacuum, defect)}}

Serialise vacuum and defect spectral results to a CSV file.  Each row
corresponds to one diffusion step~$t$.  The vacuum result provides control
measurements (unmodified Hasse graph); the defect result provides measurements
after Kuratowski contraction (Causal Prism + $K_5$ threat absorption).

\section{CSV Column Reference}

\begin{table}[H]
\centering
\begin{tabular}{@{} l l l @{}}
\toprule
\textbf{Column} & \textbf{Physics} & \textbf{Units} \\
\midrule
\code{step}       & Diffusion time $t$                          & integer steps \\
\code{P\_vac}     & Return probability on vacuum graph          & probability \\
\code{dS\_vac}    & Spectral dimension of vacuum                & dimensionless \\
\code{P\_def}     & Return probability on defect graph          & probability \\
\code{dS\_def}    & Spectral dimension with matter              & dimensionless \\
\code{P\_Gen1}    & Return probability for Generation~1         & probability \\
\code{dS\_Gen1}   & Spectral dimension, Generation~1            & dimensionless \\
\code{P\_Gen2}    & Return probability for Generation~2         & probability \\
\code{dS\_Gen2}   & Spectral dimension, Generation~2            & dimensionless \\
\code{P\_Gen3}    & Return probability for Generation~3         & probability \\
\code{dS\_Gen3}   & Spectral dimension, Generation~3            & dimensionless \\
\code{P\_Anti1}   & Return probability for antimatter           & probability \\
\code{dS\_Anti1}  & Spectral dimension of antimatter            & dimensionless \\
\code{Flux\_Attr} & Causal flux: Gen1$\to$AntiGen1 (attraction) & probability \\
\code{Flux\_Repu} & Causal flux: Gen1$\to$Gen1 (repulsion)      & probability \\
\code{Flux\_Attr\_Norm} & Normalized flux: Attr / $|\text{targets}|$ (C3) & per-node prob \\
\code{Flux\_Repu\_Norm} & Normalized flux: Repu / $|\text{targets}|$ (C3) & per-node prob \\
\code{P\_Sterile}  & Return prob.\ for sterile prisms ($N > 5$)   & probability \\
\code{dS\_Sterile} & Spectral dimension of sterile prisms (C6)    & dimensionless \\
\code{Mass\_Gen1..3} & Avg topological mass per generation         & integer (avg) \\
\code{Mass\_Anti1}   & Avg topological mass for antimatter        & integer (avg) \\
\code{dS\_vac\_std}  & Std dev of $\dS$ vacuum across realisations & dimensionless \\
\code{dS\_def\_std}  & Std dev of $\dS$ defect across realisations & dimensionless \\
\code{dS\_Gen1..3\_std} & Std dev of $\dS$ per generation          & dimensionless \\
\code{dS\_Anti1\_std}   & Std dev of $\dS$ antimatter              & dimensionless \\
\code{dS\_Sterile\_std} & Std dev of $\dS$ sterile                 & dimensionless \\
\code{Flux\_Attr\_std}  & Std dev of flux attraction                & probability \\
\code{Flux\_Repu\_std}  & Std dev of flux repulsion                & probability \\
\bottomrule
\end{tabular}
\caption{CSV output columns (32 total).}
\label{tab:csv}
\end{table}


% ███████████████████████████████████████████████████████████████████████████████
\chapter{Module Reference: \modname{memory} --- RAM-Aware Mode Selection}
% ███████████████████████████████████████████████████████████████████████████████

\begin{introduction}
  \item Automatic detection of available system RAM.
  \item Memory estimation heuristics for each simulation tier.
  \item Interactive mode selection with config file override.
\end{introduction}

\section{Functions}

\subsection{\func{load\_config(dir)} $\to$ \code{Config}}
Load configuration from a TOML file, or return defaults.  Searches for
\code{causal\_set.toml} in the given directory first, then falls back to the
current working directory.

\subsection{\func{ensure\_config(dir)}}
Write the default config file if it doesn't already exist.

\subsection{\func{recommend\_mode(n, cfg)} $\to$ \code{(ExecMode, u64, u64)}}
Recommend an execution mode based on config + available RAM vs estimated usage.
The estimate accounts for concurrent realisations (default: 2 for $N > 10$M, 4
otherwise; overridable via \code{-{}-threads}).

\subsection{\func{prompt\_mode(recommended, available, estimated, cfg)} $\to$ \code{ExecMode}}
Prompt the user to confirm or override the recommended execution mode.

\subsection{\func{estimate\_memory\_bytes(n)} $\to$ \code{u64}}
Heuristic peak memory estimate for in-memory mode:
\begin{itemize}
  \item $N \le 3{,}000$: ${\sim}N^2 \times 24$ bytes (dense eigendecomp matrices).
  \item $3{,}000 < N \le 50{,}000$: ${\sim}N \times 2{,}000$ bytes (sparse structures).
  \item $N > 50{,}000$: ${\sim}N \times 500$ bytes (OCI + CSR + defect).
\end{itemize}

\subsection{\func{estimate\_streaming\_memory\_bytes(n)} $\to$ \code{u64}}
Peak memory estimate for streaming mode: $13N + 100$\,MB.
Components: $8N$ (degree arrays) $+ 50 \cdot (N/10)$ (core CSR) $+ 100$\,MB
(buffer).  At $N = 100$M: ${\sim}1.3$\,GB per realisation.

\subsection{\func{max\_concurrent\_runs(n)} $\to$ \code{usize}}
Default concurrency heuristic: 2 for $N > 10$M, 4 otherwise.  Used as the
fallback when \code{-{}-threads} is not provided.

\subsection{\func{available\_ram\_bytes()} $\to$ \code{u64}}
Detect available system RAM in bytes via the \code{sysinfo} crate.

\section{Types}

\begin{definition}[ExecMode]
An enum: \code{InMemory} (keep all data in RAM) or \code{Streaming} (single-pass
edge analysis via \func{scan\_edges\_with\_analysis()}, ${\sim}1.3$\,GB per
realisation at $N = 100$M).
\end{definition}

\begin{definition}[Config]
User-configurable settings: \code{max\_ram\_bytes} (hard ceiling, 0 = auto),
\code{safety\_fraction} ($0.0$--$1.0$), \code{output\_dir} (default output
directory).
\end{definition}


% ███████████████████████████████████████████████████████████████████████████████
\chapter{Usage}
% ███████████████████████████████████████████████████████████████████████████████

\section{Options}

\begin{table}[H]
\centering
\begin{tabular}{@{} l l l @{}}
\toprule
\textbf{Argument} & \textbf{Description} & \textbf{Default} \\
\midrule
\code{N}            & Number of spacetime events  & 5000 \\
\code{M}            & Ensemble realisations       & 10 \\
\code{output\_dir}  & Directory for results       & current dir \\
\code{-{}-stream}         & Force streaming mode                    & --- \\
\code{-{}-inmemory}       & Force in-memory mode                    & --- \\
\code{-{}-seed} \textit{u64}       & Base seed (each realisation uses seed+i) & system clock \\
\code{-{}-threads} \textit{usize}  & Max parallel realisations               & auto (2--8) \\
\code{-{}-epsilon} \textit{f64}    & Topological error tolerance $\varepsilon$ & 0.01 \\
\code{-{}-tmax} \textit{usize}     & Maximum diffusion time                  & 15 \\
\code{-{}-eigen-cutoff} \textit{N} & Eigendecomp threshold                   & 3000 \\
\code{-{}-help}           & Show help message                       & --- \\
\bottomrule
\end{tabular}
\end{table}

\section{Configuration}

Place a \code{causal\_set.toml} in the working directory (auto-generated on
first run if absent):

\begin{verbatim}
# Hard limit on RAM (GB). 0 = auto-detect from system.
max_ram_gb = 0

# Fraction of available RAM considered safe (when auto-detecting).
safety_fraction = 0.70

# Default output directory for results and streaming edge files.
output_dir = "."
\end{verbatim}

\section{Examples}

\begin{verbatim}
# Standard in-memory run
cargo run --release -- 50000 10

# Streaming run (large N, single-pass O(N))
cargo run --release -- 1000000 5 --stream

# Reproducible run with explicit seed
cargo run --release -- 50000 10 --inmemory --seed 12345

# 16 parallel realisations on a high-RAM workstation
cargo run --release -- 10000000 20 --inmemory --threads 16

# Force eigendecomp up to N=5000 for exact comparison
cargo run --release -- 5000 10 --inmemory --eigen-cutoff 5000

# Custom walker precision (fewer walkers for exploratory runs)
cargo run --release -- 50000 10 --inmemory --epsilon 0.05 --tmax 10

# Custom output directory
cargo run --release -- 50000 10 /path/to/output
\end{verbatim}


% ███████████████████████████████████████████████████████████████████████████████
\chapter{Expected Results}
% ███████████████████████████████████████████████████████████████████████████████

\section{Spectral Dimension}

\begin{itemize}
  \item \textbf{Vacuum (global):} $\dS$ diverges toward ${\sim}15.79$ at late
        diffusion times (boundary saturation of the unstructured bulk).
  \item \textbf{Defect (with Causal Prism matter):} $\dS$ stabilises at ${\sim}8.36$
        (topological retention --- the Prism bipartite structure traps probability flow).
  \item \textbf{Coupling constant:} $\alpha_{\text{trans}} \approx 0.31$,
        consistent with the observed strong coupling at the tau mass scale.
\end{itemize}

\begin{result}
Mass is not a continuous parameter assigned to a particle, but a quantized
topological property: the integer count $N$ of intermediate nodes in the Causal
Prism.  This topological inertia stabilises the spectral dimension at a finite
value, while the vacuum diverges.
\end{result}

\section{Generation Classification}

Causal Prisms are classified by their topological signature (normalised bulk
momentum vector of poles and intermediates):

\begin{table}[H]
\centering
\begin{tabular}{@{} l l l @{}}
\toprule
\textbf{Generation} & \textbf{Physics Analogue} & \textbf{Signature} \\
\midrule
Gen1     & Electron-like  & Most abundant $[M_0, M_1, M_2, M_3]$ \\
Gen2     & Muon-like      & Second most abundant \\
Gen3     & Tau-like       & Third most abundant \\
AntiGen1 & Positron-like  & CPT conjugate of Gen1 \\
\bottomrule
\end{tabular}
\end{table}

\section{Causal Flux}

Directed walker transmission probabilities measure electromagnetic-like
interactions:
\begin{itemize}
  \item \textbf{Attraction:} Gen1 $\to$ AntiGen1 (opposite charge, positive
        flux).
  \item \textbf{Repulsion:} Gen1 $\to$ Gen1 (same charge, lower transmission).
\end{itemize}


% ███████████████████████████████████████████████████████████████████████████████
\chapter{Dependencies}
% ███████████████████████████████████████████████████████████████████████████████

\begin{table}[H]
\centering
\begin{tabular}{@{} l l @{}}
\toprule
\textbf{Crate} & \textbf{Purpose} \\
\midrule
\code{nalgebra} & Dense eigendecomposition (Tier~1, $N \le$ cutoff; default 3k) \\
\code{sprs}     & Sparse matrix algebra for $A^2$ Hasse reduction \\
\code{rand}     & Deterministic PRNG (seeded \code{StdRng}) \\
\code{rayon}    & Data-parallel iteration across walkers and nodes \\
\code{sysinfo}  & System RAM detection for automatic mode selection \\
\code{smallvec} & Stack-allocated small vectors for inner-loop optimisation \\
\bottomrule
\end{tabular}
\caption{Rust crate dependencies.}
\end{table}


% ███████████████████████████████████████████████████████████████████████████████
\chapter*{See Also}
\addcontentsline{toc}{chapter}{See Also}
% ███████████████████████████████████████████████████████████████████████████████

\begin{itemize}
  \item \emph{Modulo Synthesis Vol~I: The Geometry of Time} --- Axiom~1, causal
        sets, spectral dimension flow, the Geometry of Logic ($O(N)$ Locality
        Theorem), Law of Large Numbers.
  \item \emph{Modulo Synthesis Vol~II: The Geometry of Matter} --- Axiom~2,
        Kuratowski's Theorem, generation theorem, colour confinement, gauge
        forces, the Three Geometries.
  \item \emph{Emma's Thought Experiments} --- Pedagogic companion with
        intuitive explanations, dimensional adjustment, and locality.
\end{itemize}

\vfill
\begin{center}
{\color{main}\rule{0.6\textwidth}{1pt}}\\[1em]
{\large\bfseries\color{structure} Author}\\[4pt]
Juan Pablo Silva Alvarado\\[1em]
{\large\bfseries\color{structure} License}\\[4pt]
MIT License\\[1em]
\itshape\color{gray!60}
``The universe is strictly computable, finite, and structurally homeostatic.\\
The rest is just counting.''
\end{center}

\end{document}
